\chapter{Related Work}
\label{chap:related-work}

\chapterabstract{This chapter positions the thesis within prior work on human-scene interaction synthesis. It organizes the literature along two axes that are central to the thesis: trained versus zero-shot methods, and static 3D versus dynamic 4D interaction synthesis.}

\section{Contact Representations and Supervision for 3D HSI}
\label{sec:contact-representations}

Human-scene interaction requires more than placing a plausible human pose inside a 3D environment. A synthesized body must also be geometrically compatible with the surrounding scene: it should be supported by appropriate surfaces, avoid severe interpenetration, and make contact with scene regions that are meaningful for the intended interaction. This distinction is especially important for static 3D HSI, where the output is a single body configuration rather than a motion sequence. In such a setting, contact provides one of the main cues that separates a merely plausible pose from a physically and functionally grounded interaction.

The scene representation determines whether such contact can be enforced metrically. Indoor reconstruction datasets such as ScanNet and ScanNet++ provide RGB observations, camera calibration, and reconstructed scene geometry, allowing image-space evidence to be related to a common 3D coordinate frame \citep{dai2017scannet, yeshwanth2023scannetpp}. This calibration is essential for methods that reason jointly in image and scene space. A contact cue observed in an RGB view is only useful for 3D synthesis if it can be associated with a surface in the reconstructed scene. In this thesis, ScanNet++ provides the calibrated image, camera parameters, and scene mesh on which body-scene contact constraints are ultimately enforced.

On the human side, parametric body models provide a differentiable representation for optimization. SMPL-X is particularly suitable for human-scene interaction because it represents the full body as a posed mesh with articulated body, hand, face, shape, and global transformation parameters \citep{pavlakos2019smplx}. Its dense surface makes it possible to define contact not only at sparse joints, but over anatomically meaningful body regions. For example, grasping, standing, leaning, sitting, and pushing involve different parts of the body surface. Therefore, contact-aware HSI methods often require a mapping from semantic body parts to subsets of body vertices, rather than treating the body as an undifferentiated mesh.

Scene-aware human fitting demonstrated early that image-based human estimates are often physically invalid when evaluated in 3D. PROX extended monocular SMPL-X fitting with constraints from reconstructed scenes, using contact and interpenetration terms to reduce ambiguities in human pose estimation \citep{hassan2019prox}. Its central insight is directly relevant to synthesis: a body may explain the image evidence while still floating above the floor, penetrating furniture, or failing to touch a supporting surface. By introducing the scene as an active constraint, PROX established contact and collision avoidance as fundamental components of 3D human-scene reasoning.

Subsequent work developed contact into a learned representation for HSI. PLACE models proximity and contact between an articulated human body and a 3D scene, emphasizing that realistic bodies in scenes depend on local body-scene relationships rather than only global placement \citep{zhang2020place}. POSA further represents interaction in a body-centric way by augmenting SMPL-X vertices with contact probabilities and semantic scene labels \citep{hassan2021posa}. These methods show that contact can serve as an intermediate representation between human pose and scene affordance. They also motivate body-surface reasoning: plausible HSI depends not only on whether a person is near an object, but on which body surface is near or touching which scene surface.

Dense contact datasets and predictors further highlight the importance of body-level contact supervision. RICH introduced high-quality human-scene capture with dense vertex-level body-scene contact labels, and BSTRO uses this data to infer body-scene contact from a single image \citep{huang2022rich}. DECO extends dense 3D contact estimation to in-the-wild images by predicting vertex-level contact on the human body using both body-part and scene-context cues \citep{tripathi2023deco}. These works are important because they move beyond coarse contact categories and sparse joints. They treat contact as a dense surface property, which is more appropriate for physical interaction than global action labels alone.

Human-object interaction datasets provide another source of contact supervision. GRAB captures whole-body grasping motions with detailed body, hand, object, and contact information \citep{taheri2020grab}. BEHAVE provides multi-view RGB-D recordings with 3D human and object fits, together with contact annotations for full-body human-object interaction \citep{bhatnagar2022behave}. These datasets are valuable because they capture how the full body participates in object interaction, not only the hands. However, they are primarily object-centric. The target object is usually explicitly present and tracked, and the interaction is observed or captured rather than inferred from an open-vocabulary instruction in an arbitrary reconstructed room.

Together, these works establish that contact is a central representation for 3D HSI. They also reveal several limitations that matter for this thesis. First, many methods model contact primarily on the human side, for example by predicting which SMPL or SMPL-X vertices are in contact. This is useful for understanding the body, but it does not by itself identify the exact scene surface that should be contacted in a new reconstructed environment. Second, contact supervision is often tied to specific datasets, object categories, scenes, or captured interaction types. This limits generalization to open-vocabulary instructions and unseen scene layouts. Third, object-level grounding is often too coarse. Even when the correct object is known, the method must still determine the appropriate contact patch: a hand lifting a table should contact a plausible region of the tabletop or edge, a hand opening a washing machine should contact the door or handle region, and a foot climbing a ladder should contact a rung or step rather than the object instance as a whole.

This thesis builds on the contact-centered view of HSI, but uses contact differently. Instead of learning a contact prior from paired human-scene interaction data, the method defines body-side contact regions on SMPL-X and estimates the corresponding scene-side contact regions from visual interaction evidence. The scene-side target may be a functional part, such as a handle or control, but it may also be a pose-dependent contact patch on a larger surface, such as a tabletop, sofa seat, wall, rail, or floor. The resulting image-space masks are projected into the reconstructed scene mesh and used as explicit 3D contact targets for optimization. In this way, the method inherits the geometric accountability of contact-aware HSI while addressing a broader problem: open-vocabulary, body-part-specific scene-side contact grounding in calibrated reconstructed scenes.

\section{Static and Language-Controlled HSI Synthesis}
\label{sec:static-language-hsi}

While contact representations define how human-scene compatibility can be described, HSI synthesis methods address the generative problem: producing a human body or motion that is plausible in a given scene. In the static 3D setting, the goal is to synthesize a single body configuration that is semantically meaningful and physically feasible with respect to the environment. This requires deciding where the body should be placed, how it should be oriented, which pose it should take, and how it should relate to nearby scene surfaces. Compared with scene-agnostic pose generation, static HSI synthesis must account for both the semantic role of the scene and its metric geometry.

Early work on 3D human affordance prediction framed this problem as estimating which human poses are supported by an indoor environment. Putting Humans in a Scene predicts human affordances in 3D indoor scenes by learning what poses are plausible for a given scene, such as sitting on a chair or standing on the floor \citep{li2019putting}. The method combines semantic knowledge from 2D human poses with geometric knowledge from 3D indoor scenes, allowing it to predict semantically plausible and physically feasible human configurations. This line of work is important because it treats the scene as an active source of affordances rather than as a passive background. However, the interaction is still controlled indirectly through learned affordance distributions, not through open-ended language instructions or body-part-specific contact requirements.

Generating 3D People in Scenes without People further develops static HSI synthesis by generating full SMPL-X bodies in empty 3D scenes \citep{zhang2020generating}. The method first predicts semantically plausible body configurations conditioned on a latent scene representation and then refines the generated bodies with scene constraints to improve physical feasibility. This two-stage structure is influential: a learned model proposes a plausible human-scene configuration, while a refinement stage enforces physical constraints such as contact support and non-penetration. The method demonstrates that realistic static HSI requires both semantic compatibility and geometric correction. Nevertheless, the generated interactions are primarily governed by learned scene affordances rather than by a user-specified, open-vocabulary interaction with explicit contact intent.

Contact-aware synthesis methods such as PLACE and POSA also contribute to static HSI generation. PLACE uses local scene geometry to synthesize articulated humans that are compatible with nearby surfaces, emphasizing proximity and contact between the body and the environment \citep{zhang2020place}. POSA learns a body-centric representation of contact and scene semantics on SMPL-X, allowing plausible body placements to be generated or refined according to predicted body-scene relationships \citep{hassan2021posa}. These methods show that contact is not only useful for evaluating a generated body, but can also guide synthesis itself. However, because the contact priors are learned from available interaction data, their behavior is tied to the distribution of scenes, poses, and contact configurations seen during training.

Semantic control is introduced more explicitly in COINS, which synthesizes humans interacting with 3D scenes from action-object specifications \citep{zhao2022coins}. Instead of generating a plausible body anywhere in the scene, COINS allows the user to specify interactions such as sitting on a chair or touching a table. This is an important step from generic scene population toward controllable HSI synthesis. COINS also supports compositional interaction specifications, making it possible to combine simpler action-object relations into more complex human-scene configurations. However, the control signal remains a predefined action-object pair. This is suitable for interactions whose geometry is largely determined by the object category and action label, but it is less expressive for instructions that require fine-grained contact reasoning, such as selecting where on a large table surface the hands should touch, or which visible part of an object should be contacted.

The distinction between object-level control and contact-region control is important for this thesis. A semantic specification such as `touch table'' or `sit chair'' identifies the target object and broad action, but it does not necessarily determine the exact scene-side contact patch. Even when the relevant object part is large, such as a tabletop, sofa seat, wall, or floor, the precise contact region should still be consistent with the intended pose, body scale, orientation, and scene layout. For functional interactions, the required contact may be even more localized, such as a handle, rail, rung, button, or edge. Therefore, static HSI synthesis requires more than choosing a plausible object and body pose; it may also require grounding each relevant body part to a specific region of the scene surface.

Language-conditioned HSI extends semantic control beyond fixed action-object labels, but most existing work in this direction focuses on motion rather than a single static body configuration. HUMANISE introduced language-conditioned human motion generation in 3D scenes by aligning motion sequences with indoor scenes and language descriptions \citep{wang2022humanise}. This formulation shows that natural language can specify not only the action but also the interacting object and spatial intent. TextSceneMotion further decomposes scene-aware text-to-motion generation into target-object grounding and object-centric motion generation \citep{cen2024textscenemotion}. This decomposition is relevant because it separates the problem of understanding the language target from the problem of generating the human behavior around that target. TeSMo similarly studies text-controlled scene-aware motion generation using diffusion models, producing motions such as navigation and object interaction in 3D scenes \citep{yi2024tesmo}.

Other dynamic HSI methods broaden the scope from static placement to long-horizon scene-aware behavior. SceneDiffuser formulates scene-conditioned generation, optimization, and planning as a diffusion process applicable to tasks such as human pose generation, human motion generation, grasping, and navigation \citep{huang2023scenediffuser}. TRUMANS scales dynamic HSI supervision by capturing diverse human interactions in indoor scenes and training a model for scene-conditioned interaction sequences \citep{jiang2024trumans}. These works are important because they show how language, motion priors, scene context, and contact can be integrated for temporally evolving interactions. However, their central problem is different from the one considered in this thesis. They must model trajectories, transitions, temporal coherence, and changing contacts over time, whereas the present work focuses on one static interaction state.

Overall, static and language-controlled HSI synthesis has progressed from affordance-based human placement, to contact-aware body generation, to semantic and language-conditioned interaction control. These methods establish that plausible HSI depends on both scene semantics and physical compatibility. However, most data-driven approaches remain constrained by learned interaction distributions, predefined action-object labels, or paired scene-motion supervision. They also typically ground interaction at the level of an object, action, or motion trajectory, rather than explicitly estimating the body-part-specific scene regions that should participate in contact. This leaves two coupled challenges for zero-shot static 3D HSI: first, extracting interaction-specific 2D contact regions on the scene from visual and language evidence; and second, converting those regions into metric 3D contact targets that can guide body optimization. The next section discusses foundation-model-based methods that address parts of this problem, and motivates the need for an explicit contact-localization and geometric-grounding stage.

\section{Foundation-Model and Zero-Shot Interaction Generation}
\label{sec:foundation-zero-shot-hsi}

The methods discussed so far rely primarily on learned interaction priors, paired human-scene data, predefined action-object labels, or motion supervision. Recent foundation-model-based approaches follow a different strategy: instead of training a task-specific HSI model from paired 3D interaction data, they use pretrained language, vision-language, image generation, video generation, segmentation, and human reconstruction models as sources of interaction knowledge. These models have absorbed broad visual and semantic regularities about people, objects, actions, and everyday scenes. As a result, they can provide useful priors for interactions that are difficult to cover exhaustively with captured 3D HSI datasets.

In zero-shot HSI, foundation models are typically not used as final geometric solvers. A language model may infer the likely interaction structure from a text prompt, a vision-language model may judge whether a generated image satisfies an instruction, and an image generation model may synthesize a plausible human interacting with the scene. However, these outputs are usually expressed in language or image space. They do not directly provide a body mesh in the scene coordinate frame, nor do they automatically identify which reconstructed scene vertices should be contacted by which body parts. Therefore, zero-shot HSI requires an additional grounding step that converts foundation-model hypotheses into metric constraints.

GenZI is a key example of zero-shot static 3D HSI \citep{li2024genzi}. Given a 3D scene, an interaction description, and a specified interaction location, GenZI uses large vision-language and image inpainting models to generate plausible 2D human-scene interaction hypotheses from rendered scene views. These generated images provide visual evidence about the human pose and the intended interaction, after which a 3D human body is optimized in the scene. GenZI demonstrates that pretrained 2D models contain useful priors for 3D human-scene interaction, even without training on paired 3D HSI data. At the same time, the generated image mainly acts as an implicit interaction hypothesis. The method does not explicitly extract body-part-specific scene-side contact masks and use them as direct metric targets for optimization.

FunHSI is more directly concerned with open-vocabulary functional interaction \citep{liu2026funhsi}. It observes that many interaction prompts cannot be solved by object category alone: a person increasing room temperature, opening a window, or using a particular appliance must reason about the functional elements of the scene and the body configuration required to interact with them. FunHSI uses foundation models to identify functional scene elements, construct contact graphs, synthesize a human performing the task, and refine the 3D body through optimization for physical and functional correctness. This is closely related to the motivation of this thesis because both approaches treat open-vocabulary interaction as a contact-grounding problem rather than merely as object-level placement.

The present thesis differs in the way the scene-side contact is grounded. Instead of relying primarily on explicit functional-element identification, the proposed method first represents the interaction as body-part-to-scene contact requirements through a Scene Interaction Graph. The scene node may be a target object or the floor, and the precise contact patch is inferred later from visual evidence. This formulation covers functional contacts, such as touching a handle or control, but it also covers broader pose-dependent contacts, such as choosing where on a tabletop, wall, sofa seat, treadmill rail, or floor the body should touch. Thus, the problem is not only functional object-part discovery. It is broader: interaction-specific scene-side contact localization for each required body part.

This distinction is important because the difficult step is not only projecting 2D contact evidence into 3D. Projection is necessary, but it assumes that an appropriate image-space contact region has already been identified. In many interactions, obtaining that 2D region is itself ambiguous. A generated person may visually suggest that a hand is lifting a table, but the method must still decide which visible patch of the tabletop or edge corresponds to the hand contact. A person may be shown reaching toward a washing machine, but the correct hand contact should lie near the door or handle region rather than an arbitrary visible surface. A broad object mask is not sufficient in either case. The contact region must be consistent with the instruction, the generated body pose, the body-part identity, the visible object geometry, and the original calibrated scene image.

A related development outside HSI is the use of foundation image models as components of agentic 3D reasoning systems. WorldAgents investigates whether 2D foundation image models can be used as agents for 3D world synthesis, using a multi-agent architecture with a vision-language director, an image generator, and a verifier that evaluates generated frames in both image and reconstruction space \citep{erkoc2026worldagents}. Although its goal is 3D world generation rather than human-scene interaction, it supports a broader observation that is relevant here: foundation image models become more useful for 3D tasks when embedded in iterative generate-and-verify loops rather than used as one-shot image synthesizers.

This motivates the agentic contact-localization stage used in this thesis. The generated human frame is treated as a visual reference for pose, orientation, laterality, and likely contact layout, while the original scene image remains the authoritative canvas. An image model proposes colored overlays for the required scene-side contact regions; the colored pixels are extracted and recomposited onto the unedited original image; and a vision-language evaluator checks whether the resulting composite correctly marks the contact regions implied by the generated interaction. If the contact masks are incorrect, a correction instruction is issued and the process is repeated. The output of this loop is not a final image, but a set of verified body-part-specific 2D contact regions that can be projected to the reconstructed scene mesh.

Several zero-shot human-object interaction methods are related to this direction. DreamHOI uses text-to-image diffusion priors for zero-shot 3D human-object interaction synthesis with a given human and object, optimizing the human articulation so that the body appears to interact realistically with the object \citep{zhu2024dreamhoi}. InterFusion addresses text-driven 3D human-object interaction generation by using text-derived human pose priors and a two-stage generation process to reduce the difficulty of directly producing a full interaction scene from text \citep{dai2024interfusion}. InteractAnything uses large language models and diffusion-based object affordance parsing for zero-shot 3D human-object interaction, including extraction of contact points for open-set objects \citep{zhang2025interactanything}. These works show that foundation models can support interaction synthesis beyond closed object and action datasets.

However, object-centric HOI differs from full-scene HSI. In HOI, the target object is usually explicit, isolated, or provided as an input asset. In HSI, the object is embedded in a larger scene with surrounding geometry, occlusions, support surfaces, and other objects. A body may need to contact both the target object and the floor, avoid nearby obstacles, and remain consistent with the original reconstructed environment. Therefore, scene-side contact grounding in HSI must account not only for object affordance, but also for camera calibration, visible scene geometry, mesh projection, and body-scene collision constraints.

Part-level affordance reasoning in HOI is nevertheless relevant. HOI-PAGE introduces Part Affordance Graphs for zero-shot human-object interaction generation, representing how human body parts engage with object parts and using these relations to guide 4D HOI synthesis \citep{li2025hoipage}. This supports an important trend: realistic interaction generation increasingly requires structured relations between body parts and scene or object parts, rather than only global text prompts or whole-object proximity. The Scene Interaction Graph used in this thesis follows a related motivation, but it is intentionally simpler on the scene side. It specifies which body regions should contact the target object or floor, while the exact scene-side contact region is estimated visually from the generated interaction and original scene image.

Foundation-model priors have also been explored for dynamic HSI. ZeroHSI uses video generation models as motion priors for zero-shot 4D human-scene interaction, reconstructing temporally evolving human-scene interactions through neural human rendering and differentiable optimization \citep{li2024zerohsi}. GenHSI similarly uses pretrained video diffusion models for controllable human-scene interaction video generation, decomposing the process into script writing, 3D-aware keyframe previsualization, and animation \citep{li2025genhsi}. These works show that generated videos can provide priors over human motion and interaction dynamics. Their goal, however, is temporal visual synthesis or 4D interaction generation, whereas this thesis focuses on a single static body configuration with explicit metric contact constraints.

Overall, zero-shot interaction generation methods show that foundation models provide useful semantic, visual, and motion priors for HSI and HOI. They reduce dependence on paired 3D interaction datasets and enable broader language control. The remaining challenge is to make these priors geometrically accountable. For static 3D HSI, this means first extracting interaction-specific 2D contact regions on the original scene image, and then projecting those regions to the reconstructed mesh as scene-side targets for body optimization. The proposed method addresses this gap by using foundation models not only to generate a plausible human interaction image, but also to localize and verify body-part-specific contact regions before enforcing them through SMPL-X optimization.

\section{Research Gap and Positioning}
\label{sec:gap-positioning}

Prior work has established several important ingredients for human-scene interaction synthesis. Contact-aware 3D HSI methods show that plausible bodies in scenes require contact and collision reasoning. Static synthesis methods show that humans can be generated in 3D environments using learned scene affordances, body-scene compatibility priors, or semantic action-object control. More recent zero-shot methods show that foundation models can provide useful visual and semantic priors for open-vocabulary interactions. However, these directions do not fully address the problem considered in this thesis: synthesizing a static 3D human in a reconstructed scene from an open-vocabulary instruction while grounding each required body contact to a specific scene-side region.

The key missing step is interaction-specific scene-side contact grounding. Object-level grounding is often too coarse, because knowing the target object does not determine where the body should touch it. Part-level grounding is also not always sufficient. Some interactions involve small functional parts, such as handles, controls, rungs, or buttons, but others involve large object parts or broad support surfaces. For example, lifting a table may require hand contact on a plausible region of the tabletop or table edge; sitting on a sofa requires selecting an appropriate patch on the seat or back support; and leaning on a wall requires choosing a contact region consistent with the body pose and orientation. Thus, the desired contact target is not simply an object instance, and sometimes not even an entire semantic object part, but a localized body-part-specific region of the visible scene surface.

Foundation-model-based methods help by producing plausible visual interaction hypotheses, but these hypotheses are not directly usable as 3D constraints. A generated image may suggest that a hand touches an object, but it does not automatically provide a reliable 2D mask for the contacted scene region, nor a corresponding set of 3D mesh vertices. Therefore, zero-shot static HSI requires two linked operations: first, extracting the correct 2D contact regions on the original scene image; and second, projecting those regions into the reconstructed scene as metric contact targets.

This thesis addresses this gap through an agentic contact-grounding pipeline. A Scene Interaction Graph first specifies which body regions should contact the target object or floor. A generated human frame then provides a visual hypothesis for pose, orientation, laterality, and approximate contact layout. Instead of accepting this generated image as the final output, the method uses it to guide contact localization on the original calibrated scene image. Colored contact overlays are generated, extracted, recomposited onto the unchanged original image, and verified by a vision-language evaluator. This closed-loop process produces body-part-specific 2D contact masks while reducing the effect of scene drift or incorrect one-shot annotations.

The accepted masks are then projected to the reconstructed scene mesh and used as explicit scene-side targets for SMPL-X optimization. The final body is initialized from monocular human reconstruction and refined using contact attraction, penetration avoidance, and initialization priors. This distinguishes the thesis from prior work in three ways: it does not require paired 3D HSI training data, it is not limited to predefined action-object labels, and it treats generated visual evidence as an intermediate source of contact constraints rather than as a final visual result. The contribution is therefore the conversion of open-vocabulary visual interaction evidence into metrically grounded body-part-specific contact targets for static 3D human-scene interaction synthesis.
